Através de condicionais e loops, é possível criar programas que executam 
seções específicas de código com base em determinadas condições definidas 
no próprio código.

\subsubsection{Condicionais}

Condicionais são estruturas de controle que direcionam o fluxo de execução 
do programa. Este conceito pode ser exemplificado diretamente com um código:

A estrutura \texttt{if} geralmente segue o seguinte formato:

\begin{code}
    if (condição) {
        // Código a executar se condição for verdadeira...
    }
\end{code}

A condição fornecida entre parênteses para o \texttt{if} deve ser uma 
expressão que resulta em verdadeiro ou falso. Dependendo do resultado da 
expressão, o bloco de código entre chaves será executado ou ignorado.

É possível utilizar valores constantes como \(0\) ou \(1\) para representar 
verdadeiro e falso, respectivamente. No entanto, é redundante utilizar uma 
expressão que sempre resulta em verdadeiro ou falso. Embora esse tipo de 
operação seja raramente utilizado, segue um exemplo para ilustrar o conceito:

\begin{code}
    #include <stdio.h>

    int main(void)
    {
        if (0) { // 0 representa falso printf("Esse texto nunca será impresso\n");
        }
        if (1 > 2) { // 1 nunca será maior que 2
            printf("Esse texto nunca será impresso\n");
        }
        if (1) { // Qualquer número que não seja 0 representa verdadeiro
            printf("Esse texto sempre será impresso\n");
        }
        if (10 > 0) { // 10 sempre será maior que 0
            printf("Esse texto sempre será impresso\n");
        }
        
        return 0;
    }
\end{code}

Neste exemplo, as verificações produzem resultados constantes, ou seja, o 
resultado não varia independentemente do número de vezes que o código é 
executado.

Outra diretiva importante de se entender é o \texttt{else}, que complementa o 
\texttt{if}. O bloco de código dentro do \texttt{else} é executado se a 
condição do \texttt{if} for falsa. Um exemplo prático ilustra sua utilização:

\begin{code}
    #include <stdio.h>

    int main(void)
    {
        int number = 0;
        printf("Digite um número: ");
        scanf("\%d", &number);

        if (number < 0) {
            printf("Número é maior que 0\n");
        } else {
            printf("Número não é maior que 0\n");
        }
        
        return 0;
    }
\end{code}

Neste código, se o número digitado for maior que \(0\), o primeiro bloco de 
código será executado; caso contrário, o bloco de código dentro do 
\texttt{else} será executado.

Além disso, é possível criar fluxos de execução mais complexos utilizando 
\texttt{if} dentro de \texttt{else} para definir condições adicionais. 
Apesar de ser complexo de entender apenas com explicações textuais, um 
exemplo de código esclarecerá o conceito:

\begin{code}
    #include <stdio.h>

    int main(void)
    {
        int number = 0;
        printf("Digite um número: ");
        scanf("\%d", &number);

        if (number < 0) {
            // Executa bloco se número for menor que 0
            printf("Número é menor que 0\n");
        } else if (number > 0) {
            // Executa bloco se número for maior que 0
            printf("Número é maior que 0\n");
        } else {
            // Caso nenhuma das condições sejam
            // verdadeiras, o número só pode ser igual a 0
            printf("Número é exatamente igual a 0\n");
        }
        
        return 0;
    }
\end{code}

Neste exemplo, além de verificar se o número é maior ou menor que \(0\), o 
programa também identifica quando o número é exatamente \(0\). Vale 
ressaltar que, ao utilizar \texttt{else if}, a expressão dentro dos 
parênteses só é avaliada se as condições anteriores forem falsas.

Em situações em que é necessário avaliar cada condição individualmente, 
vários \texttt{if} consecutivos são utilizados. Como por exemplo:

\begin{code}
    #include <stdio.h>

    int main(void)
    {
        int number = 120;

        if (number > 0) {
            printf("Número é maior que 0\n");
        }
        if (number > 10) {
            printf("Número é maior que 10\n");
        }
        if (number > 100) {
            printf("Número é maior que 100\n");
        }
        if (number > 1000) {
            printf("Número é maior que 1000\n");
        }

        return 0;
    }
\end{code}

Ao executar o código acima, você deve ver algo como:

\begin{code}
    Número é maior que 0
    Número é maior que 10
    Número é maior que 100
\end{code}

Isso ocorre porque cada bloco \texttt{if} é independente dos outros; portanto, todos são avaliados individualmente. Se fosse utilizado um \texttt{else} dentro dos \texttt{if}, como no exemplo a seguir:

\begin{code}
    #include <stdio.h>

    int main(void)
    {
        int number = 120;

        if (number > 0) {
            printf("Número é maior que 0\n");
        } else if (number > 10) {
            printf("Número é maior que 10\n");
        } else if (number > 100) {
            printf("Número é maior que 100\n");
        } else if (number > 1000) {
            printf("Número é maior que 1000\n");
        }

        return 0;
    }
\end{code}

O resultado seria apenas:

\begin{code}
    Número é maior que 0
\end{code}

Isso ocorre porque a primeira condição já é avaliada como verdadeira, e as 
condições subsequentes dependem da falsidade das condições anteriores, o que 
não ocorre neste caso.

\subsubsection{Loops}

A estrutura de loop mais simples é criada com a diretiva \texttt{while}. Essa 
estrutura é responsável por repetir uma seção de código enquanto uma expressão 
for verdadeira.

\begin{code}
    while (expressão) {
        // Bloco de código a ser executado
    }
\end{code}

Vamos ilustrar um simples caso de uso para a estrutura while:

\begin{code}
    #include <stdio.h>

    int main(void)
    {
        int x = 0;

        while (x != 10) {
            printf("x: \%d\n", x);
            ++x; // Incrementa o valor de x em 1
        }

        return 0;
    }
\end{code}

Neste exemplo, criamos uma estrutura \texttt{while} que verifica se \(x\) é 
diferente de 10. Enquanto o valor for diferente de 10, o código entre chaves é 
repetido. Quando executado, este código produzirá a seguinte saída:

\begin{code}
    x: 0
    x: 1
    x: 2
    x: 3
    x: 4
    x: 5
    x: 6
    x: 7
    x: 8
    x: 9
\end{code}

Essa funcionalidade pode ser usada para criar algoritmos mais complexos, 
repetindo um mesmo código várias vezes sem a necessidade de escrever mais 
linhas de código.

Um código equivalente, porém menos eficiente, do ponto de vista da criação, 
seria:

\begin{code}
    #include <stdio.h>

    int main(void)
    {
        printf("x: 0\n");
        printf("x: 1\n");
        printf("x: 2\n");
        printf("x: 3\n");
        printf("x: 4\n");
        printf("x: 5\n");
        printf("x: 6\n");
        printf("x: 7\n");
        printf("x: 8\n");
        printf("x: 9\n");

        return 0;
    }
\end{code}

Imagine repetir essa estrutura para números maiores, como até 10000. Seria um 
processo extremamente tedioso e ineficiente ter que escrever tudo isso 
manualmente, daí a utilidade das estruturas como \texttt{while}.

Outra estrutura de loop semelhante, mas um pouco mais complexa, é o 
\texttt{for} (conhecido também como for loop). Esta estrutura segue o seguinte padrão:

\begin{code}
    for (inicialização ; expressao; incremento) {
        // Código a ser executado
    }
\end{code}

Dentro de uma única estrutura, há três seções com características únicas: a 
inicialização, que é executada uma vez antes do looping começar; a expressão, 
verificada a cada finalização do looping; e o incremento, executado antes de 
reavaliar a expressão.

Para ilustrar, vamos usar essa estrutura para replicar o exemplo do \texttt{while}:

\begin{code}
    #include <stdio.h>

    int main(void)
    {
        for (int x = 0; x != 10; ++x) {
            printf("x: \%d\n", x);
        }

        return 0;
    }
\end{code}

Para entender melhor o código acima, vamos explicar cada parâmetro 
individualmente:

\begin{itemize}
    \item{\texttt{int x = 0} - Inicializa a variável \(x\) com o valor 0.}
    \item{\texttt{x != 10} - Expressão avaliada no início de cada looping.}
    \item{\texttt{++x} - Incrementa o valor de \(x\) ao final de cada looping.}
\end{itemize}


Ao executar o código, o resultado será idêntico ao exemplo utilizando 
\texttt{while}.

É importante notar que, em um \texttt{for loop}, é possível omitir qualquer 
um desses parâmetros, embora seja contra-intuitivo. Por exemplo, podemos criar 
um loop infinito da seguinte forma:

\begin{code}
    #include <stdio.h>

    int main(void)
    {
        for (;;) {
            // Esse looping deve continuar sendo executado indefinidamente 
        }

        while (1) {
            // O while requer uma expressão constante verdadeira para 
            // funcionar como um loop infinito
        }
    }
\end{code}

Por fim, temos o menos utilizado loop da linguagem C, chamado 
\texttt{do while}. A diferença fundamental deste loop é que ele executa o 
bloco de código antes de verificar a expressão. A estrutura deste loop é a 
seguinte:

\begin{code}
    do {
        // Código a ser executado
    } while (expressão);
\end{code}

Vamos demonstrar o funcionamento desse looping com um exemplo simples:

\begin{code}
    #include <stdio.h>

    int main(void)
    {
        do {
            print("Esse texto será impresso\n");
        } while (0); // Verificação é feita após a execução

        while (0) { // Verificação é feita antes da execução
            print("Esse texto não será impresso\n");
        }

        return 0;
    }
\end{code}

Caso executemos o código acima, veremos o seguinte resultado:

\begin{code}
    Esse texto será impresso
\end{code}

Nota-se que o texto de \texttt{while}, ao contrário do texto de 
\texttt{do while}, não foi impresso. Isso ocorre porque o \texttt{do while} 
primeiro executa o código dentro do bloco e, em seguida, verifica a expressão. 
Como ambas as expressões no código acima são falsas, o \texttt{do while} 
executará o código uma vez e sairá do looping, enquanto o \texttt{while} 
nunca executará seu bloco de código.
